\subsection{Classes of unitary representations}

Lemma 9. --- Let E be a Hilbert space over K. Let F ⊂ E be a dense vector subspace of E. Then the Hilbert dimension of E is less than or equal to the dimension of F.

If E has finite Hilbert dimension, this is equal to the dimension of E. Suppose that E has infinite Hilbert dimension. The subspace F is then of infinite dimension. Let B be an orthonormal basis of E and let B' be a basis of F. For every \(x \in B\) , there exists an element \(f(x) \in B'\) such that \(\langle x \mid f(x) \rangle \neq 0\) , since otherwise we would have \(x \in F^{\circ} = \{0\}\) . One thus defines a map \(f: B \to B'\) .

For every \(y \in B'\) , the set \(f^{-1}(y)\) is contained in the set of \(x \in B\) such that \(\langle x | y \rangle \neq 0\) and is therefore countable (EVT, V, p. 21, prop. 4). By E, III, p. 50, prop. 4, one obtains \(\text{Card(B)} = \text{Card}(f(B)) \leqslant \text{Card(B')}\) .

Denote by \(Is_{G}(\pi_{1},\pi_{2})\) the relation

“G is a topological group and \(\pi_{1}\) and \(\pi_{2}\) are

unitary representations of G in Hilbert spaces over K that are

isomorphic.” With respect to \(\pi_{1}\) and \(\pi_{2}\) , this is an equivalence relation. For any unitary representation \(\pi\) of G in a Hilbert space over K, we shall denote by \(\mathrm{cl}(\pi)\) the equivalence class of \(\pi\) (cf. E, II, p. 47); it is therefore a unitary representation of G isomorphic to \(\pi\) ; one says that \(\mathrm{cl}(\pi)\) is the class of \(\pi\) . Unitary representations \(\pi_{1}\) and \(\pi_{2}\) in Hilbert spaces over K are isomorphic if and only if \(\mathrm{cl}(\pi_{1}) = \mathrm{cl}(\pi_{2})\) .

Let G be a topological group. Let c be a cardinal. The relation

“λ is a class of unitary representation of G in

a complex Hilbert space of Hilbert dimension \(\leqslant\mathfrak{c}\)”

is collectivizing in \(\lambda\) (E, II, p. 3). Indeed, every Hilbert space over K of dimension \(\leqslant \mathfrak{c}\) is isometrically isomorphic to a Hilbert subspace of \(\ell^2 (\mathfrak{c})\) (EVT, V, p. 23, cor. 2), and the assertion then follows from E, II, p. 47.

Let \(\pi\) be an irreducible unitary representation of G in a Hilbert space E. Let \(x\) be a nonzero element of E. Since \(\pi\) is irreducible, the vector \(x\) is a cyclic vector of \(\pi\) , which implies that the Hilbert dimension of E is \(\leqslant\) Card(G) (Lemma 9, applied to the dense subspace generated by the elements \(\pi(g)x\) for \(g \in G\) ). The classes of irreducible unitary representations of G in a Hilbert space over K therefore belong to the set of classes of unitary representations of G in a Hilbert space over K of Hilbert dimension \(\leqslant\) Card(G); consequently, they form a set.

DEFINITION 8. --- We denote by \(\widehat{\mathbf{G}}\) the set of classes of irreducible unitary representations of \(\mathbf{G}\) in a complex Hilbert space. We say that \(\widehat{\mathbf{G}}\) is the unitary dual of \(\mathbf{G}\) .

For every irreducible unitary representation \(\pi\) of G in a complex Hilbert space, we therefore have \(\operatorname{cl}(\pi) \in \widehat{\mathbf{G}}\) .

Remarks. --- 1) Suppose that G is a commutative group. By Corollary 7 of V, p. 390, every irreducible unitary representation of G in a complex Hilbert space has dimension 1. If G is locally compact, the set \(\widehat{G}\) is identified with the set of unitary characters of G (def. 1 of II, p. 201), and the notation \(\widehat{G}\) is therefore compatible with that introduced in Def. 2 of II, p. 201.

\begin{enumerate}
\def\labelenumi{\arabic{enumi})}
\setcounter{enumi}{1}
\item
  If G is finite, the set \(\widehat{\mathbf{G}}\) is in bijection with the set of classes of \(\mathbf{C}[\mathrm{G}]\) -simple modules (A, VIII, p. 47), which is also denoted \(\widehat{\mathbf{G}}\) in A, VIII, p. 396.
\item
  For \(\pi \in \widehat{\mathbf{G}}\) , we shall identify \(\overline{\pi}\) with the class in \(\widehat{\mathbf{G}}\) of the conjugate representation of \(\pi\) .
\item
  If \(\pi\) is an irreducible unitary representation of G in a finite-dimensional complex Hilbert space, its character \(\chi_{\pi}\) depends only on the class of \(\pi\) . One can therefore speak of the set of characters of finite-dimensional complex irreducible unitary representations of G.
\end{enumerate}

